\pdfmapfile{+cm.map}
\pdfmapfile{+cmextra.map}
\pdfmapfile{+symbols.map}
\documentclass[11pt]{article}
\usepackage[margin=0.85in]{geometry}
\usepackage{amsmath,amssymb,amsthm,longtable,booktabs,array,hyperref}
\hypersetup{colorlinks=true,linkcolor=blue,urlcolor=blue,pdftitle={Sparse carry encodings for the consistency of ZFC},pdfauthor={OpenAI-assisted research for Vladimir Reshetnikov}}
\newtheorem{theorem}{Theorem}
\newtheorem{lemma}{Lemma}
\providecommand{\mathscr}[1]{\mathcal{#1}}
\newcommand{\Z}{\mathbb Z}\newcommand{\N}{\mathbb N}
\newcommand{\Con}{\operatorname{Con}}
\title{Sparse carry encodings for the consistency of ZFC}
\author{OpenAI-assisted research for Vladimir Reshetnikov and ProveIt}
\date{23 October 2026}
\begin{document}\maketitle
\begin{abstract}
The purpose of this supplement is an explicit, reproducible improvement to fixed Diophantine presentations of $\Con(\mathrm{ZFC})$, not a proof of consistency. The central change replaces exponentially spaced exponent positions by a certified four-sum-distinct sequence. A second change specializes the carry compiler to a fixed polynomial, eliminating an unused input digit and an entire transfer test. A third packs an entire finite tile payload into one complement equality. We distinguish the universal mathematical arguments, exact finite certificates, inherited compiler results, and external theorem dependencies. No end-to-end proof-assistant certificate or global minimum is claimed.
\end{abstract}
\section{Question, domains, and review boundary}
The repository snapshot is
\texttt{22e75dd6b87b1b199bb08c778d7cf075ae5aae08}.
It contains seven wave-five packages numbered 31--37. Package 37 is the preceding delivery in this conversation. The comparison must distinguish positive integers, natural numbers $\N=\{0,1,\ldots\}$, and unrestricted integers. In every construction, named circuit registers abbreviate expressions; they are not further existential unknowns. All coefficient leaves are fixed integers with finite recipes, and no recipe branches on the truth of consistency.

The desired assertion has the form
\[
\mathrm{ZFC}\vdash
\bigl(\neg\exists\mathbf x\in D^r\; F(\mathbf x)=0\bigr)
\ \longleftrightarrow\ \Con(\mathrm{ZFC}).
\]
The proof proceeds by showing that a solution is equivalent to a finite syntactic refutation of ZFC. It does not decide either side. The literal 215-state machine, its logical calibration, and the ordinary native-window/Pell compiler are inherited dependencies. The new arithmetic adapters are proved here in ordinary mathematics and checked by reproducible finite certificates, not claimed to have been rebuilt in Lean or Rocq.

\section{A four-sum-distinct exponent set}
\label{sec:b4}
Let $M=2^{20}-1=1048575$. Work in the field represented by
\[
\mathbb F_2[T]/(T^{20}+T^3+1).
\]
The certificate verifies that $\theta=T$ has an orbit of length $M$ containing every nonzero bit vector. Thus every nonzero element is invertible and $\theta$ generates the multiplicative group. The 32 elements fixed by $x\mapsto x^{32}$ form a subfield. For each such element $a$, take the discrete logarithm of $\theta+a$.

These 32 logarithms have distinct sums of four terms modulo $M$, allowing repetitions. Indeed equality of such sums gives equality of two monic products $\prod_{i=1}^4(\theta+a_i)$. Subtracting the products gives a polynomial of degree at most three over $\mathbb F_{32}$ vanishing at $\theta$. The extension has degree four, so that polynomial is zero and unique factorization gives equal multisets.

Translation and multiplication by a unit modulo $M$ preserve this property. Translate one member to zero, leaving 31 positive representatives $w_i$. Padding a multiset with zero converts uniqueness of four-term sums into uniqueness of all sums $\sum\alpha_iw_i$ with $|\alpha|\le4$. Our deterministic affine-orbit search gives
\[
\begin{gathered}
(w_1,\ldots,w_{31})=(51609,55430,68596,78030,89463,90233,\\
110498,110508,128615,136498,211819,241519,305379,312706,\\
314312,369115,417382,429683,435112,450230,462314,469608,\\
490199,501539,507151,524944,550792,559943,573626,607799,610895).
\end{gathered}
\]
The actual machine-readable list is authoritative; the displayed list is for inspection. All $\binom{35}{4}=52360$ four-term multisets are also enumerated and their residues checked independently. The multiplier is 208418 and the subsequent translation is 836756, applied after subtracting 535448 from the raw logarithms modulo $M$. No optimality among all four-sum-distinct sets is asserted.

Put $t=610895$ and $L=4t=2443580$. The former exponential spacing used $5^{31}$ or $5^{32}$. It is uniqueness of low-degree sums, not that particular spacing, that is needed for coefficient extraction.

\section{A fixed-polynomial carry lemma}
Let $P\in\Z[X_1,\ldots,X_{31}]$ have degree at most four and $P(0)>0$. Write
\[
 c_\alpha=(4-|\alpha|)!\prod_i\alpha_i!\,[X^\alpha]P,
 \qquad D(B)=\sum_\alpha c_\alpha B^{L-\sum_i\alpha_i w_i}.
\]
Choose fixed powers of two $z>\max_\alpha|c_\alpha|$ and $a$ such that
\[
 a>\max\{2z\,32^4,2048z,8\},\qquad b\ge2,\qquad B=ab^4.
\]
For digit witnesses $0\le n_i<b$, put $g=\sum_i n_iB^{w_i}$. The coefficient of $X^L$ in
\[
 (1+\sum_i n_iX^{w_i})^4\sum_\alpha c_\alpha X^{L-\alpha\cdot w}
\]
is exactly $24P(\mathbf n)$. This follows from the four-sum property and the multinomial coefficient. Every coefficient has absolute value less than $z(1+\sum n_i)^4<B/2$.

If $b$ is a power of two, define
\[
 J_c=2(1+g)^4D(B)+B\sum_{j=0}^{8t}B^j,
 \qquad T_c=(B-2)B^L.
\]
Then binary addition of $J_c$ and $T_c$ has no carry exactly when $P(\mathbf n)=0$. To see this, divide both operands by two: the shifted coefficients lie strictly between zero and $B$, and the mask tests whether the coefficient at position $L$ is precisely $B/2$.

This argument uses a signed sparse coefficient polynomial. It needs neither the dense positive coefficient background nor an enormous fixed index that transports its value from one base to another. Detailed two-sided bounds, the witness-digit mask, central-binomial encoding, and the ratio interface are supplied in the accompanying proof files.

\section{Results and comparison}
The following numbers describe explicit fixed equations, not unspecified representatives of recursively enumerable sets. Arithmetic operations count binary addition, subtraction and multiplication, including multiplication by numerals; fixed coefficient construction is excluded and specified separately. Domain is part of every claim.
\begin{center}\small
\begin{tabular}{@{}lrrr@{}}\toprule
Construction & Unknowns & Exact degree & Operations\\\midrule
Sun sparse equation & $10\Z$ & $69\,554\,088\,184$ & 1407\\
Jones sparse equation & $9\N$ & $23\,321\,536\,264$ & 1383\\
Positive translate & $9\Z_{>0}$ & $23\,321\,536\,264$ & 1378\\
Jones merged guard & $13\Z$ & $226\,334\,238\,192$ & 1409\\
Jones four guards & $14\Z$ & $226\,295\,141\,096$ & 1420\\
Unit inner separator & $12\Z_{>0}$ & 901 & 112\\
Shorter carrier & $12\Z_{>0}$ & 905 & 110\\
Rebound $P_{67}$ & $15\Z_{>0}$ & 155 & 67\\\bottomrule
\end{tabular}
\end{center}
The thirteen-integer construction does not use Sun's 2026 Lucas theorem; it is a genuinely different dependency choice, not a competitor to ten integers on variable count alone. Its fourteen-integer predecessor has slightly smaller degree and is retained for that reason. The twelve-positive equations also give three-equation systems of maximum degree 892 or 896, with 106 or 104 operations in their shifted presentations. The new coefficient tuple rebinds the seven-equation, sixteen-positive-variable system at degree 24 and 59 operations; those operation and variable counts are unchanged.

\paragraph{Comparison with the reviewed snapshot.}
Package 31 had ten integers at degree $126496\cdot5^{31}+27064$ and 1442 operations. Packages 34/37 had nine-variable degree $92324\cdot5^{32}+296$, with package 37 costing 1505 over $\N$ and 1500 over positives; older package 27 had 1473 operations at a larger degree. The new nine-positive circuit beats even that historical cost. Package 34's lowest twelve-positive degree was 1022 (124 operations), with intermediate tradeoffs 1023/119, 1025/116, 1033/114 and 1041/111. The new 901/112 and 905/110 forms are meaningful improvements, but the old 108-operation degree-1081 form and 102-operation shifted system remain relevant on cost alone.

The smallest bound coefficient base in the reviewed collection was $2^{1012115}$, with 506664 raw payload bytes. An optional unbound tuple in package 33 had 506644 bytes. The complement-coded tuple below has base $2^{590069}$ and 295388 raw bytes, improving both comparisons. It is bound explicitly to $P_{67}$ and the 59-operation system only. The new low-variable equations use the complete older package-31 tuple, not a mixture of its ports with the new tuple.

There is no single justified ``simplest'' form across all these metrics. No sub-ten unrestricted-variable or sub-nine natural-variable result is claimed. We found no validated thirty-variable parent quartic and no operation count below 67/59. The concrete contribution is instead a change in the exponent coding, a merged signed guard, a unit inner separator, and a new coefficient decoder.

\section{A ten-integer equation with radically smaller degree}
\label{sec:sun-b4}
The new source compiler preserves the same ten unrestricted integer
unknowns and the same fixed native coefficient tuple as the pinned
\texttt{31-sun10} equation. Its exact degree is
\[
 \boxed{69\,554\,088\,184=113856\cdot610895+27064},
\]
versus $589\,042\,901\,992\,797\,851\,562\,527\,064$ previously, a reduction
by approximately $8.47\times10^{15}$. Its emitted evaluator uses 1,407
binary operations rather than 1,442: 484 additions, 25 subtractions and
898 multiplications. The multiplication count increases from 851, so
the total-operation improvement should not be described as improving
every circuit metric. Both circuits have 79 fixed leaves.

\paragraph{The changed source.}
Let $P=\sum_{j=1}^{22}r_j(n+1)^2$ be the pinned quartic in 31 nonnegative
unknowns. For $P=\sum p_\alpha n^\alpha$, put
$c_\alpha=(4-|\alpha|)!\prod_j\alpha_j!p_\alpha$.
Use the certified $B_4$ witness positions $w_1,\ldots,w_{31}$ with
$t=\max w_j=610895$, and put $L_c=4t$.
Keep $\beta_S=2^{4269882}$, which is distinct from the native compiler
port named \texttt{beta}. Define
\begin{align*}
 b&=1+3f,&\mathcal B&=\beta_Sb^4,&\mathscr C&=b\mathcal B^t,\\
 \ell&=\sum_j\mathcal B^{w_j},&
 D_c&=\sum_\alpha c_\alpha\mathcal B^{L_c-\alpha\cdot w},\\
 N_0&=\mathcal B^{t+1},&N_1&=4\mathcal B^{8t+1},&N&=N_0N_1,\\
 M_c&=N_0-1-(b-1)\ell,&
 T_c&=M_c+(\mathcal B-2)\mathcal B^{L_c}N_0,\\
 J_c&=2(1+g)^4D_c+\sum_{i=0}^{8t}\mathcal B^{i+1},&S_c&=g+J_cN_0,\\
 R_c&=(S_c+T_c+1)N+T_c+1,&X&=(N-1)R_c,&Y&=N^2.
\end{align*}
All exponents are fixed integers. There is no dummy source input slot.
The companion proof establishes the necessary bounds for arbitrary
$0\le g<2\mathscr C$, before assuming that $g$ decodes to a tuple.
Specifically, if $f\ne0$ and $b$ is square, then
$0<D_c$, $0<J_c<2\mathcal B^{8t+1}<N_1$ and $0\le S_c,T_c<N$.
Moreover $3\mid X$, $X>3b$, and $Y>\max\{b,256\}$.
For square powers of two $b$, binary carry decoding yields
\[
 Y\mid\binom{2X}{X}
 \quad\Longleftrightarrow\quad
 g=\sum_jn_j\mathcal B^{w_j},\quad 0\le n_j<b,\quad P(n)=0
\]
within the prescribed $g$ interval. Arbitrarily large $b$ supply the
positive direction. This replaces the full source theorem application
by a proved sparse specialization while preserving its entire interface.

\paragraph{The complete equation.}
The unknowns are $(f,g,h,k,l,w,x,y,z,m)\in\mathbb Z^{10}$.
Continue the preceding source polynomials by
\begin{align*}
 L&=lY,&U&=2LX,&V&=4gwY,&W&=bw,\\
 A&=U(V+1),&B&=2X+1,&C&=B+(A-2)h,\\
 P_n&=U^2V,&K&=X+1+k(P_n-2),\\
 S&=2A-5,&T&=3WC-2(W^2-1),&R_2&=(P_n^2-4)K^2+4.
\end{align*}
\[
 \Delta=8\mathscr C^3gK^2-g^2\bigl(32(C-KL)^2\mathscr C^3+g^2K^2\bigr).
\]
\begin{align*}
 O&=f^2l^2\Delta,&\xi&=zx,\\
 D&=(A^2-4)C^2+4,&E&=C^2D\xi,\\
 F&=4(A^2-4)E^2+1,&G&=2(A+2)(A-2)^2E^2-1,\\
 H&=C+(B-C)F+16EFy,&I&=(G^2-1)H^2+1.
\end{align*}
Set
\begin{align*}
 r&=T+Sy,&\Sigma&=\xi S^5,&
 \Theta&=(r^3-rS^2-S^3)(r^2+23S^2),\\
 a_0&=b\bigl(1+(3O-4)b^2z^2\bigr),&a_1&=R_2,&a_2&=DFI,\\
 v_0&=1,&v_1&=a_0,&v_2&=a_0(P_n^2+K^2+3),\\
 Z&=\Sigma m-\Theta,&A_i&=\Sigma^2v_i^2a_i\quad(0\le i\le2),\\
 V_n&=Z^2+A_2-A_0-A_1,&W_n&=V_n^2+4Z^2A_2-4A_0A_1.
\end{align*}
Then the new equation is
\[
 \boxed{F_{B_4}=W_n^2-16V_n^2Z^2A_2=0.}
\]
These are abbreviations, not additional unknowns. The bridge and finalizer
are unchanged from the pinned construction; only the source coding changed.
The source-interface theorem therefore gives the same equivalence between
integer zeros of $F_{B_4}$ and nonnegative zeros of $P$.
Together with the inherited calibrated compiler endpoint,
\[
 \mathrm{ZFC}\vdash
 [\neg\exists(f,g,h,k,l,w,x,y,z,m)\in\mathbb Z^{10}\ F_{B_4}=0]
 \longleftrightarrow\operatorname{Con}(\mathrm{ZFC}).
\]

\paragraph{What was proved and what was checked.}
The new ingredient is the sparse source-interface proof, including its
unconditional size bounds. It relies on the explicitly certified $B_4$
positions and standard binary carry identities. Sun 2026 Theorem~4.1 and
Propositions~6.1--6.2, and the repository's compiler calibration, remain
external mathematical inputs. No new universal unknown-count theorem,
global optimality, coefficient-size reduction, or end-to-end Lean/Rocq
formalization is claimed.

The exact degree is attained by nonzero leading diagonal residues 95,
909 and 668382 modulo 101, 1009 and 1000003. The finite audit also verifies
all 402 coefficient recipes, all fixed payload hashes, all 52,360 padded
four-sums, and the formal 35-term final norm identity. A separate literal
transcription matches 120 modular assignments and 5,400 named values;
these last tests are regressions rather than universal equivalence proofs.
The full proof, circuit and checks are under \texttt{sun/}.

\section{Nine natural unknowns: positive coding and a near-square majorant}
\label{sec:nine-b4}
The sparse carry construction also gives a single equation in nine
nonnegative unknowns. Its exact degree is
\[
  38176t+8744=23\,321\,536\,264,\qquad t=610895,
\]
with 1383 binary arithmetic operations. The positive-integer translate uses
1378 operations and has the same degree. These are ordinary mathematical
reductions to the inherited literal compiler and Jones ratio theorems,
not new kernel-checked consistency results or variable-minimality claims.

The fixed source is the 31-witness parent with SHA256
\texttt{1c336ba0b3c9} (full digest in the artifact manifest),
bound to the complete 121-by-8801, 943-tile tuple in
\path{sun/sources/data/coefficients_binding.json}.
Its ports $(\beta,K_0,M_C,M_F)$ have bit lengths
$(1064921,1070004,1064921,1064922)$, and $M_F$ is already shifted.
For $P_0(\mathbf n)=\sum r_j(\mathbf n+1)^2$, the 402 weighted
coefficients satisfy $|c_\alpha|+1<z=2^{4269855}$ and
$c_0=24P_0(0)>0$. Put $a=2^{4269877}$, so
$a>\max(2z\,32^4,2048z,8)$.

\paragraph{Front end and unconditional positivity.}
The unknowns are $(\epsilon,g_0,h_0,s_0,w_0,\phi,i_0,j_0,n)\in\mathbb N^9$.
Set $g=g_0+1$, $h=h_0+1$, $s=s_0+1$, $w=w_0+1$, $i=i_0+1$,
$j=j_0+1$. For the certified positive $B_4$ weights $w_\ell$, put
$L=4t$, $m_c=8t+1$, $e_c=9t+2$, and
\begin{align*}
 b&=2+\epsilon,& B&=ab^4,& q_0&=B^{t+1},& \ell&=\sum_kB^{w_k},\\
 D_c&=\sum_\alpha c_\alpha B^{L-\alpha\cdot w},&
 M_c&=q_0-1-(b-1)\ell,\\
 J_c&=2(1+g)^4D_c+B\sum_{k=0}^{m_c-1}B^k,&
 S_c&=g+q_0J_c,\\
 T_c&=M_c+q_0(B-2)B^L,& N&=2B^{e_c},\\
 R&=S_c(N^2-N)+(T_c+1)(N^2-1),&G_{\rm cap}&=32bB^t.
\end{align*}
The sparse carry proof applies once the ratio margin forces
$1\leq g<4bB^t$. The quartic radix is valid because
$|2(1+g)^4D_c|<1024zb^4B^{8t}<B^{8t+1}/2$, while every
coefficient of the decoded product has magnitude
$<z32^4b^4<B/2$. Also $B>4b$, and $B$ is a power of two whenever $b$ is. Crucially, positivity holds even without that bound:
the unique leading coefficient of $D_c$ is $c_0\geq1$, hence
\[
 D_c(B)\geq B^L-z\sum_{k<L}B^k>0\quad(B-1>z).
\]
The digit mask $M_c$ is nonnegative; thus $J_c,S_c>0$, $T_c\geq0$,
and $R\geq N^2-1\geq N\geq b$. The old negative-$R$ corrections
are therefore unnecessary for this fixed source.

\paragraph{Ratio registers.}
All the following are circuit registers, not additional unknowns:
\begin{gather*}
 U=N^2w,\quad Y=N^2s,\quad M=RY,\quad P=2M^2U,\\
 K=R+1+h(P-1),\quad A=M(U+1),\\
 B_1=2R+1,\quad W_1=bw,\quad C=B_1+W_1+\phi,\quad J=4A-5,\\
 Q=3W_1C-2(W_1^2-1),\quad D=(A^2-1)C^2+1,\\
 E=iJDC^2,\quad F=(A^2-1)E^2+1,\\
 G=A+F(CD+1-A),\quad H=B_1+jC,\quad I=(G^2-1)H^2+1,\\
 a_0=(P^2-1)K^2+1,\quad a_1=64DFI,\quad u=JF,\\
 v=QF+J(H-C),\quad
 \mu=K^2G_{\rm cap}-4(C-KY)^2G_{\rm cap}-8gK^2.
\end{gather*}
They satisfy $A\geq2$, $J\geq3$, $D\geq A^2$ and
$C,D,E,F,G,H,I>0$. Also $C>W_1$ and $H-C>0$, so $Q,u,v>0$ and
\[
 0<\frac vu=\frac QJ+\frac{H-C}{F}\leq Q+H-C=:M_{\rm div}.
\]

\paragraph{The new structural majorant.}
Write $t_0=iJC$, so $E=t_0DC$. Then
\begin{align*}
 DF&=t_0^2D^3(D-1)+D,\\
 S_0&=t_0(8D^2-4D-1),\\
 S_0^2-64DF&=t_0^2(8D+1)-64D\geq8D+9>0.
\end{align*}
Here $t_0^2\geq9$. This exact identity is a polynomial near-square
approximation, not an asymptotic argument. Since $S_0>0$ and
$I\leq G^2H^2$, the choices
\[
 V_0=PK,\qquad V_1=S_0GH
\]
obey $\sqrt{a_k}<V_k$; here $K>1$ makes the first bound strict. The factor 64 does not alter the
square predicate: $64x$ is a square over $\mathbb Z$ exactly when $x$ is.

\paragraph{The final equation and both directions.}
Define
\[
 \Omega=u(2\mu-1),\qquad
 \Theta=un+v-\Omega(M_{\rm div}+V_0V_1).
\]
The polynomial is
\[
 \boxed{\ F_9=(\Theta^2-\Omega^2(a_0+V_0^2a_1))^2
              -4\Omega^4V_0^2a_0a_1=0.\ }
\]
For completeness, its positive-divisor finalizer is as follows. If
$u,v>0$, $M\geq v/u$, and $\sqrt{|a_k|}<V_k$ with $V_k\geq2$,
then this polynomial has a natural $n$ exactly when both $a_k$ are
integer squares, $u\mid v$, and $\mu>0$.
Indeed, factor it over roots $r_k^2=a_k$ as
$\prod_{\eta,\xi=\pm1}[\Theta-\Omega(\eta r_0+\xi V_0r_1)]$.
A vanishing factor makes $s=r_0+V_0r_1=\Theta/\Omega$ rational.
For $s\ne0$, $2sr_0=s^2+a_0-V_0^2a_1$ makes both roots integer.
For $s=0$, $|r_1|=|r_0|/V_0<1$ forces $a_1=a_0=0$.
Thus no square-class independence is assumed. Now integrality gives
$|r_k|\leq V_k-1$. In all cases
\[
 Z=M+V_0V_1+r_0+V_0r_1\geq M+1>0,
 \qquad n+v/u=(2\mu-1)Z.
\]
The left side is positive, so $\mu>0$, and integrality forces $u\mid v$.
Conversely, integer roots and these conditions give the integer
$n=(2\mu-1)Z-v/u\geq1$. Eliminating $n$ therefore recovers exactly
the original four Jones ratio conditions. The sparse carry theorem,
fully bound quartic, and literal compiler endpoint give
\[
 \mathsf{ZFC}\vdash
 [\neg\exists\mathbf x\in\mathbb N^9\ F_9(\mathbf x)=0]
 \ \longleftrightarrow\ \operatorname{Con}(\mathsf{ZFC}).
\]

\paragraph{Exact degree, rather than a syntactic upper bound.}
The leading terms of $V_1^2$ and $a_1$ cancel. Resolve this exactly by
\begin{align*}
 \Delta_V=V_1^2-a_1
  &=[t_0^2(8D+1)-64D]G^2H^2
       +64DF(H^2-1),\\
 E_0&=un+v-\Omega M_{\rm div},\\
 \Theta^2-\Omega^2(a_0+V_0^2a_1)
  &=E_0^2-2E_0\Omega V_0V_1
       +\Omega^2V_0^2\Delta_V-\Omega^2a_0.
\end{align*}
On diagonal specialization, $\Delta_V$ has degree $9072t+2072$
and leader $2^{243}a^{2268t+504}$. The last display's third term uniquely dominates
with degree $19088t+4372$ and leader $2^{523}a^{4772t+1060}$; the final norm's
subtracted term has lower degree $31264t+7164$. Hence
$\deg F_9=38176t+8744$ with diagonal leader $2^{1046}a^{9544t+2120}$.
The generated telescoping certificates audit these exact cancellations.

Replacing each natural coordinate by a positive coordinate minus one is
a bijection preserving degree. In particular $b=1+\epsilon_+$ and
$C=2R+bw+\phi_+$ permit the reported five-operation saving.
The verifier reconstructs all 402 coefficient recipes, authenticates all
four coefficient payloads, checks all 52,360 four-multiset sums, expands
the new identities exactly, and checks 240 independent modular instances
and both domain translations. These finite checks validate the artifacts;
they do not replace the universal proof or the imported compiler theorems.

\section{Thirteen integers using only the Jones ratio theorem}

There is a second, independent route to a small signed equation which does
not use Sun's 2026 reduction.  It combines the fixed-polynomial carry code
above with the literal Jones ratio construction and a new guard that
simultaneously certifies an interval and an accurate rational approximation.
The output has
\[
 \boxed{13\text{ unrestricted integers},\quad 1409\text{ binary operations},
 \quad\deg F_{13}=226334238192.}
\]
More precisely its degree is $370496t+84272$, where $t=610895$.
This is not a variable-count improvement over the companion ten-integer
construction.  It is an explicit smaller Jones-based alternative with no
2026 Lucas-reduction dependency.  The fourteen-integer variant has $1420$
operations and slightly smaller degree $226295141096$; both are retained.

\subsection*{Exact specification}
The full division-free recipe is
\path{signed_audit/generated/signed13_b4.json}; its source coefficients
are the same coherent $31$-variable quartic tuple as the companion
construction.  Set $a=2^{4269877}$ and use precisely the variables
\[
 (\epsilon,g,h,s,w,\phi,i,j,z_1,z_2,z_3,v,u)\in\mathbb Z^{13}.
\]
Put $b=2\epsilon^2$, $\omega=w^2$, and use the fixed-polynomial B4 front end
with $B=ab^4$, $t=610895$, $L=4t$, and $N=2B^{9t+2}$.  It produces $R$ and
\(c=bB^t\).  Retain the literal Jones ratio registers
\[
\begin{gathered}
 U=N^2\omega,\quad Y=N^2s,\quad M=RY,\quad P=2M^2U,\quad
 K=R+1+h(P-1),\quad A=M(U+1),\\
 B_1=2R+1,\quad W_1=b\omega,\quad C=B_1+W_1+\phi,
 \quad J=4A-5,\quad Q=3W_1C-2(W_1^2-1),\\
 D=(A^2-1)C^2+1,\quad E=iJDC^2,\quad F=(A^2-1)E^2+1,\\
 G=A+F(CD+1-A),\quad H=B_1+jC,\quad I=(G^2-1)H^2+1,\\
 a_0=(P^2-1)K^2+1,\qquad a_1=DFI.
\end{gathered}
\]
The new combined guard is
\[
 \boxed{\Delta=g(c-g)K^2-c^2(C-KY)^2>0.}
\]
Only $s\ge1$, $\phi\ge0$, and $\Delta>0$ remain to be encoded.  Set
\[
 d_1=3s-1,\quad d_2=3\phi+2,\quad d_3=3\Delta-1,
 \quad C_0=1,\quad C_r=C_{r-1}(1+4d_rC_{r-1}z_r^2),
 \quad C_*=C_3.
\]
Let $D_0=JF$, $T_0=QF+J(H-C)$, $Z=\epsilon wi z_1z_2z_3$, and
\[
\begin{gathered}
 S=ZD_0^5,\quad r=T_0+D_0v,\quad
 T=(r^3-D_0^2r-D_0^3)(r^2+23D_0^2),\\
 w_1=C_*,\quad w_2=C_*(P^2+K^2+2),\quad X=Su-T,\\
 A_0=S^2C_*,\quad A_1=S^2w_1^2a_0,\quad A_2=S^2w_2^2a_1,\\
 V=X^2+A_2-A_0-A_1,\quad W=V^2+4X^2A_2-4A_0A_1,\\
 \boxed{F_{13}=W^2-16V^2X^2A_2=0.}
\end{gathered}
\]
All displayed registers are polynomial abbreviations, not additional
unknowns.  The three-factor Pell stack, quintic intersective merger and
three-root norm have the full signed-domain proofs in
\path{signed_audit/SIGNED_B4_PROOF.md}; the new guard proof is in
\path{signed_audit/THIRTEEN_GUARD_PROOF.md}.

\subsection*{Why the guard is sound}
Since $c>0$, $\Delta>0$ forces $K\ne0$ and $0<g<c$.  The exact identity
\[
 c^2\bigl(K^2-4(C-KY)^2\bigr)
   =4\Delta+(c-2g)^2K^2
\]
gives the Jones approximation $4(C-KY)^2<K^2$.  Positivity of $C,Y$
then forces $K>0$; the square condition on $a_0$ excludes $h\le0$ by the
usual gap between consecutive squares.  Reflection of $i$ and a Pell
orbit modulo $CF$ restore positive $i,j$ without changing the conditions.
To invoke the unchanged \texttt{Nine.ratio\_sound} theorem literally, supply
its free ceiling parameter, in the proof only, as $G_{\rm aux}=8gK^2+1$.
The positive integer $K^2-4(C-KY)^2$ is at least one, so its original
margin is positive.  Thus the theorem gives the required power-of-two
base and central-binomial divisibility.  The recovered $g<c$ supplies
the carry-code bounds.

\subsection*{Why the canonical witness satisfies the stronger guard}
The exact source \texttt{Ratio.exists\_B25core} chooses $b=2^k$ and
\[
 \omega=2^{2R+1-k},\quad
 U=\frac{N^2}{b}2^{2R+1},\quad
 \rho=\frac{(U+1)^{2R}}{U^R}=Y+f.
\]
Expanding the binomial, the fractional tail is
\[
 f=\sum_{j=0}^{R-1}\binom{2R}{j}U^{j-R}
   <\frac{2^{2R}}U=\frac{b}{2N^2}<1.
\]
The other terms are the integer $Y$, and $N^2\mid Y$ as in the original
construction.  The literal \texttt{ratio\_core} estimate, with
$A=RY(U+1)$, yields
\[
 \left|\frac CK-\rho\right|
 \le\frac{2(Y+f)}{Y(U+1)}<\frac4U.
\]
Consequently, since $R\ge8$,
\[
 \left|\frac CK-Y\right|
 <\frac b{N^2}\left(\frac12+\frac4{2^{2R+1}}\right)
 <\frac b{N^2}.
\]
The direct witness code has $1\le g<c$, while $c,b\le N$.  Therefore
\[
 \frac{c^2(C-KY)^2}{K^2}<\frac{c^2b^2}{N^4}\le1
 \le g(c-g),
\]
which proves $\Delta>0$.  Neither $Y$ nor the base is changed by the later
product-Pell extension.  Choosing $k=1+2r$ sufficiently large makes
$b=2(2^r)^2$ and $\omega$ a positive square.  The three stack witnesses,
merger witness and norm root now complete the integer solution.

Thus ordinary mathematics formalizable in ZFC proves that $F_{13}$ has an
integer zero exactly when the source quartic does, hence exactly when
there is a finite ZFC refutation.  The accompanying source, B4, polynomial,
modular and exact-degree checks are diagnostics, not a new end-to-end
Lean certificate or a proof of consistency.

\section{A twelve-positive-variable polynomial of degree 901}
\label{sec:positive901}
The inner radical-separating multiplier can be removed, rather than merely
made cheaper.  The resulting single equation has exact degree $901$ and
$112$ binary arithmetic operations.  Its three-equation presentation has
maximum degree $892$ and $107$ operations, or $106$ after the inherited
positive transport shift.  A shorter carrier gives degree $905$ in $110$
operations (shifted system: degree $896$, $104$ operations).
Every binary $+,-,\times$, including multiplication by a numeral, is charged;
comparisons in systems are free.  The degree-$901$ form improves the pinned
$1022/124$ form in both degree and cost.  These are ordinary mathematical
reductions formalizable in ZFC, not consistency proofs or new proof-assistant
certificates.

\paragraph{Fixed instance and explicit definition.}
Select the complete source31 Wang943 $121$-by-$8801$ tuple:
$\beta=2^{1064921}-1$, $K_0=K$, $M_C=M_C$, $M_F=M_{F,\mathrm{source}}$,
with the \emph{shifted} mask.  The included coefficient binding gives the
immutable source adapter and all four materialized payload hashes.  The twelve positive
unknowns are $J,F,\alpha,z,h,i,T,s,w,\eta,\zeta,t$; all following letters
are computed polynomial abbreviations:
\begin{align*}
 q&=\beta J+1,& X&=wq,&Y&=sq^3,&E&=XY,\\
 k&=\eta+\zeta,&c&=kY+\eta,&a&=E+Y,&H&=4a+3,\\
 \Delta&=a^2+H,&\ell&=EkY,&A&=1+\ell(\ell+k),&G&=1+\Delta c^2,\\
 b_0&=X+ac,&I&=ic^2,&D&=1+\Delta I^2,&Q&=(\Delta I)^2,\\
 d&=Q-1,&v&=cT,&b&=c+RD,&e&=Q(Dv^2+b^2)-1,\\
 B_*&=2Qvb,&K_*&=IT+1,&L&=1+HK_*,&W&=b_0+H(t+b+\delta),
\end{align*}
where
\begin{align*}
R&=((q-1)(q-F)+\alpha)(q^2-1)+(M_C+qM_F)J,\\
N_k&=k-hE-R,\qquad
N_t=(K_0+w)(q-F-\alpha)+q-F-z(q-1),\\
\delta&=(a+2)c-\ell\quad\text{(full carrier)},\qquad
\delta=ac-\ell\quad\text{(short carrier)}.
\end{align*}
The integer polynomial $N$ is the eight-factor conjugate product
\[
 N=\prod_{g=\pm\sqrt G,\,\tau=\pm\sqrt A,\,f=\pm\sqrt D}
 \bigl[d(W-Lg-H\tau)^2-H^2e+H^2B_*f\bigr].
\]
Formal radical signs specify a polynomial, not extra unknowns or independent
Galois automorphisms.  The equation and system are
\[
 N_kN_t(1+N)-1=0,\qquad\text{or}\qquad N_k=1,\quad N_t=1,\quad N=0.
\]
For an explicit division-free evaluator put
\begin{align*}
 u&=dH^2A,&p&=H^2QDv^2,&r&=H^2Qb^2,&j&=4pr,\\
 e_H&=p+r-H^2,&\mathcal T&=2d(W^2+L^2G),&\mathcal K&=[d(W^2-L^2G)]^2,\\
 a_*&=\mathcal T-2(u+e_H),&b_*&=(u-e_H)^2-j-\mathcal K,\\
 v_*&=16uj-a_*(a_*\mathcal T+2b_*),&u_*&=b_*^2-a_*^2\mathcal K,\\
 N&=\mathcal K v_*^2-u_*(\mathcal T v_*-u_*).
\end{align*}
Eliminating $\tau,f$ gives $(x^2-2(u+e_H)x+(u-e_H)^2-j)^2-16ujx$;
reducing modulo $x^2-\mathcal T x+\mathcal K$ proves the formula.
Also $N\equiv(u_*-d(W^2+L^2G)v_*)^2\pmod4$.  Hence a product root
forces $N=0$ and $N_k=N_t=\epsilon\in\{1,-1\}$.

\paragraph{The two new recovery lemmas.}
The unchanged weak compiler interface gives
\[
 q\ge16,\quad Y\ge q^3,\quad E=XY\ge q^4,
 \quad 0<R<E,\quad k>E,\quad kY<c<2kY,\quad a+2<2E.
\]
For estimates let $g=\sqrt G,f=\sqrt D,\tau=\sqrt A$ be positive real roots.
Direct expansion gives $1<A<\Delta c^2\le\Delta I=\sqrt Q$ and $g>\tau$.
Moreover $g<2Ec$, $f<2EI$, $I\ge c^2$ and $c>EY$ imply
\begin{align*}
 b+cf&<5E^3I^2+2EI^2<6E^3I^2<\tfrac12E^4I^2<d,\\
 Ig-cf&=\frac{I^2-c^2}{Ig+cf}
       >\frac{I}{8Ec}=\frac{ic}{8E}>\frac Y8>1.
\end{align*}
The second line uses the exact identity $(Ig)^2-(cf)^2=I^2-c^2$.
For
\[
 C(x)=d(x^2+A)-e,\quad
 F(x)=[C(x)^2+DB_*^2-4d^2x^2A]^2-4C(x)^2DB_*^2,
\]
every complex root has the coherent form
$x=\tau'+y$, $(\tau')^2=A$, $(f')^2=D$,
$dy^2=Q(vf'-b)^2-1$.
Writing $Z=b+vf>1$ gives
\[
 |y|<Z(1+1/d),\qquad Z/d<T,
 \qquad |y|<b+ITg,\qquad |x|<b+(IT+1)g.
\]
Here the first bound follows by squaring, the second from $b+cf<d$,
and the third from $T(Ig-cf)>T$.
At a norm zero, a negative main root would instead give
\[
 x=(W+Lg)/H=b+t+\delta+b_0/H+(K_*+1/H)g>b+K_*g,
\]
a contradiction.  Irrational quadratic conjugation would force that negative
branch as well.  Thus the selected $g$ is a positive integer and
$x=(W-Lg)/H$ is rational.  This argument covers imaginary auxiliary roots.

The crucial inner recovery no longer needs a large multiplier.
If $C(x)=0$ at a root of $F$, substitution of $x^2=e/d-A$ and
$B_*=2Qvb$, followed by exact division by $4$, gives
\[
 d^2A^2-deA+Q^2v^2b^2D=0.
\]
Reducing modulo $Q$ forces $Q\mid A(A-1)$, impossible because
$0<A(A-1)<Q$.  Likewise $F(0)=0$ forces $Q\mid(A-1)^2$, also impossible.
Consequently the justified rational formulas
\[
 f'= -\frac{C(x)^2+DB_*^2-4d^2x^2A}{2C(x)B_*},\qquad
 \tau'=\frac{C(x)+B_*f'}{2dx},\qquad y=x-\tau'
\]
recover integer square roots $f',\tau'$.  The inherited rational Pell
saturation lemma applied to $Q=(\Delta I)^2$ makes $y$ integral.
Reduction modulo $H$ gives $\gamma=(g-b_0)/H\in\mathbb Z$, and
$g>(a+1)c$, $c>X$ give $\gamma>0$.  The normalized signed-quotient
compiler interface now reconstructs the accepting blank K215 computation.

\paragraph{Completeness and positivity are constructive.}
For a canonical positive compiler witness, $V=vf-b$ is a nonzero integer,
so $y^2=V^2+(V^2-1)/d\ge V^2$ and $y\ge b-vf$.
Both corrections satisfy $0<\delta<g$, since $\ell/c>E/2>2$.
Set
\[
 t=\gamma+(IT+1)g+\tau+y-b-\delta
   >\gamma+\tau+ITg-vf>0.
\]
Then $W=Lg+H(\tau+y)$, giving $N=0$ and $N_k=N_t=1$.
No new Pell congruence, enlarged index, or divisibility restriction is used.
The inherited substitution $z_{\rm old}=z+1$ remains positive after compiler
reconstruction and removes one system gate.  Composing with the existing
K215 calibration and negating proves in ZFC that absence of positive
solutions is equivalent to $\operatorname{Con}(\mathrm{ZFC})$.

\paragraph{Exact degree and verification.}
Scale all twelve unknowns by a formal parameter.  Put
$q_0=\beta J$, $k_0=\eta+\zeta$, $Y_0=sq_0^3$,
$E_0=wq_0Y_0$, $c_0=k_0Y_0$, $I_0=ic_0^2$, $H_0=4E_0$,
$\ell_0=E_0c_0$, $d_0=E_0^4I_0^2$,
$B=q_0^3(q_0-F)E_0^2I_0^2$, and $V_0=I_0T\ell_0$.
Four conjugate factors have degree $119$; cancellation of the matching
$f,g$ signs leaves degrees $107,107,105,97$ for the full carrier.
Their product has nonzero leading form
\[
 -2^{14}H_0^{16}d_0^8B^8V_0^4\ell_0^2k_0Y_0^2
 \left(2t+k_0-\frac{iTc_0}{E_0}\right),
\]
of degree $892$.  The apparent denominator cancels against the prefactors;
independent occurrence of $t$ proves nonvanishing for every fixed $\beta>0$.
For the short carrier the final factor instead has degree $101$ and the
leading form is
$2^{16}H_0^{16}d_0^8B^8V_0^4\ell_0^2k_0Y_0^2c_0$, of degree $896$.
The nonzero leading forms of $N_k,N_t$ add $7+2$ degrees.
The full proof supplies every cancellation and denominator check.
The reproducible checker independently multiplies the formal radical factors,
checks signed and degenerate fixtures, and verifies these complete polynomial
ray degrees and analytic leading coefficients over three finite fields.
Finite diagnostics do not substitute for the unbounded compiler theorem.

\section{Complement-packed native clauses}
\label{sec:complement-native}
This section gives an independent structural reduction of the complete
coefficient tuple at the inherited 67-operation polynomial and
59-operation system frontiers.  Their arithmetic circuits, degrees, and
positive-variable domains do not change.  The old tuple in package 36 had
$B=2^{1012115}$ and occupied 506664 raw bytes.  The new tuple has
\[
 \rho=2^{67},\qquad L=8807,\qquad d=590069,\qquad B=\rho^L=2^{590069},
\]
and its four materialized ports occupy 295388 raw bytes, a reduction of
$41.6994\%$.  This changes the integers themselves, rather than merely
their file compression.

\subsection{One congruence replaces the parity fields}
\begin{lemma}[Aggregate complement decoder]
Let $N<2^a$, let $0<v_i<2^m$ for $0\le i<N$, and let all $x_i,y_j$
be Boolean.  Define
\[
 n=\sum_i x_i,\quad W=\sum_{j=0}^{m-1}2^j y_j,\quad
 S=\sum_i\bigl(1+2^a(2^m-v_i)\bigr)x_i+2^aW,
 \quad \mu=2^{a+m}-2.
\]
Then $S\mathbin{\mathrm{AND}}\mu=0$ if and only if either all selectors
and copies vanish, or exactly one selector $i$ is present and $W=v_i$.
\end{lemma}
\begin{proof}
The low $a$ bits of $S$ equal $n$, since $0\le n\le N<2^a$.
Their mask is $2^a-2$, so $n\in\{0,1\}$.  This extraction precedes,
and does not assume, any claim about internal payload carries.
If $n=0$, the remaining payload is $W<2^m$, whose low $m$ bits must
vanish, so $W=0$.  If selector $i$ is present, the payload is
$2^m-v_i+W\in(0,2^{m+1})$.  Its low $m$ bits vanish exactly when it
equals $2^m$, hence exactly when $W=v_i$.  Conversely valid cells have
$S=0$ or $S=2^{a+m}+1$, both disjoint from $\mu$.
\end{proof}
The interval $W<2^m$ uses one distinct Boolean support position for each
copy bit.  It does not apply to unconstrained nonnegative copy values.
The native low mask supplies these Boolean conditions; both dummy
positions have clause coefficient zero.

Retain the exact Wang943 language and its ordering.  For a tile with
west, north, east, south labels $(w,n,e,s)$, put
\[
 v_i=w+2^8n+2^{17}e+2^{25}s+127\,2^{34}.
\]
The lower 34 bits encode the labels.  Positions $34,35,36,37$ are
the four occupancy bits, and positions $38,39,40$ are the three anchors.
These last seven bits are one for every occupied tile.  Use $a=10,m=41$,
so the selector coefficient and mask are
\[
 c_i=1+1024(2^{41}-v_i),\qquad \mu=2^{51}-2.
\]
Moving the seven constant bits to the top gives $c_i\le 1+2^{44}$.
This additionally improves the uniform carry bound.

The 38 native edge-copy positions are unchanged.  Their payload positions
are specified by
\[
\begin{array}{c|c|c}
\text{slot}&\text{native positions}&\text{payload bit positions}\\\hline
W&1023+j,\ 0\le j<9&0,\ldots,7,34\\
N&1013+j,\ 0\le j<10&8,\ldots,16,35\\
E&1003+j,\ 0\le j<9&17,\ldots,24,36\\
S&983+j,\ 0\le j<10&25,\ldots,33,37
\end{array}
\]
The anchors $1095,2195,3296$ carry payload positions $38,39,40$.
Let $E$ be these 41 positions together with selectors $0,\ldots,942$
and dummies $1063,1064$; thus $|E|=986$.  At the copy position for bit
$j$ put $c_e=2^{10+j}$, and at each dummy put $c_e=0$.  The lemma now
proves exactly the same local valid cells as the old 41 separate
radix-four parity fields.

\subsection{The full coefficient recipe}
All powers below have fixed finite exponents.  Define
\begin{align*}
 D_C={}&\rho^{30}+\rho^{62}+\rho^{1100}+\rho^{2201}
       +\sum_{e\in E}c_e(\rho^{4393-e}+\rho^{5495-e}),\\
 D_R={}&\rho^{42},\\
 Z={}&(2^{65}-1)\sum_{e=8792}^{8804}\rho^e
                  +(2^{19}-1)\rho^{8805},\\
 M_C={}&B-1-\sum_{e\in E\setminus\{1\}}\rho^e-2\rho^{1063},\\
 M_{F,\mathrm{native}}={}&4+\mu(\rho^{4393}+\rho^{5495})
       +\sum_{j=0}^9\rho^{1013+j}+\sum_{j=0}^8\rho^{1065+j}
       +\rho^{2195}+\rho^{3296}+Z.
\end{align*}
The four arithmetic ports, all of which must be replaced together, are
\[
 \boxed{\beta=B-1,\quad K_0=D_C+BD_R,\quad M_C,\quad
 M_F=M_{F,\mathrm{native}}+B-1.}
\]
The fourth port is shifted; substituting the native high mask is incorrect.

The old support geometry is preserved.  The two clause targets are
$4393,5495$; the two distinct tested temporal anchors remain $2195,3296$
with shared current source $1095$.  Finite support arithmetic verifies
\[
 ((4393-E)\bmod L)\cap((5495-E)\bmod L)=\varnothing,
 \qquad
 ((2195-E)\bmod L)\cap((3296-E)\bmod L)=\{0\}.
\]
All 66062 current/spatial/temporal native-basis identities are verified
independently.  The aggregate lemma at the clean band of the distinguished
Start cell forces its shared anchor, and hence both tested temporal
anchors, to one.  The inherited rotation argument therefore still excludes
binary misalignment and then nonzero whole-radix cell misalignment.
After alignment, overlap parity recovers the original occupied edge
codes and the same Wang matching constraints.

For the inherited conservative uniform bound one now has
\[
 \mathcal M=988\left(2\sum_{e\in E}c_e+6\right)
 =29613572795083482608,
 \qquad
 4\mathcal M+8=118454291180333930440<2^{67}.
\]
Also $2\mu+4<2^{67}$ and $2(2^{65}-1)+4<2^{67}$.  These bounds
preserve the compiler's carry, positive-slack, and nonaligned-rotation
estimates.  The new internal payload carries are intentional and are
covered by the lemma, not by an assumption of separate fields.

The largest unshifted product exponent is 8791.  Each padded test, at a
position from 8792 through 8805, therefore constrains a zero in an aligned
cell.  There are 864 such tests.  Consequently
\[
 \operatorname{pc}(M_C)=d-986,\quad
 \operatorname{pc}(M_{F,\mathrm{native}})=986,
 \quad M_C\equiv2\pmod4,\quad
 M_{F,\mathrm{native}}\equiv4\pmod8.
\]
The original native mask population balance and low congruences remain exact.

\subsection{The new actual-rank controller}
For completeness use the explicitly certified prime
\[
 \lambda=581197,\qquad k=\lambda-1=581196,
 \qquad d=67\cdot8807.
\]
The multiplicative orders of $B$ modulo $67,8807,\lambda$ are respectively
$66,4403,581196$.  Their required divisibility conditions hold, and
\[
 B^k\bmod\lambda^2=182964883980\ne1,
 \qquad B^k\equiv1\pmod\lambda.
\]
For $H_0=1+BD_C+B^2(D_R+1)$, the residues
\[
\begin{array}{c|rrr}
p&B&B^2-1&H_0\\\hline
67&46&38&41\\
8807&456&5374&5193\\
581197&348469&573553&76889
\end{array}
\]
prove the necessary unit conditions.  Choose sufficiently large independent
spatial and temporal dimensions $w,h$ that are powers of $\lambda$,
and put $N=wh$, $q=B^N$.  Switching a safe low dummy bit at position $j$
changes the actual packed rank by
\[
 \delta R=-\Gamma B^j,\qquad
 \Gamma=\rho^{1063}(q^2-1)[1+q(D_C+BD_R+B^w)].
\]
The certificates imply $\gcd(\Gamma,dN)=1$.  With $T=N/\lambda$,
the powers $B^{kj}$ run through all residues $1+d\lambda z$ modulo
$dN$.  The original Boolean consecutive-orbit subset argument, using
$\{B^{2+kj}:0\le j<T\}\cup\{B^3\}$, therefore adjusts the actual
rank to $R\equiv dw\pmod{dN}$.  Taking $N>d\lambda^2$ and $h>\lambda$
keeps all chosen dummy cells and their needed translates below $N$ and
away from the origin.  The complete argument is reproduced in
\texttt{native/PROOF.md}; it is not a test of only synthetic tableaux.

Both $d$ and these canonical dimensions are odd.  The unchanged population
and congruence identities retain the source square/odd-scale promises.
Soundness, unlike completeness, places no prime-power restriction on the
dimensions of an arbitrary recovered tableau.

\subsection{The short explicit system}
Use sixteen strictly positive unknowns
\[
 J,F,\alpha,z,f,h,i,o,j,s,w,\tau,\eta,\zeta,y,\gamma.
\]
Define the following polynomial abbreviations, not additional unknowns:
\begin{gather*}
q=\beta J+1,\quad C=q-F-\alpha,\quad X=wq,\quad Y=sq^3,\\
E=XY,\quad k=\eta+\zeta,\quad c=kY+\eta,\quad a=E+Y,\\
H=4a+3,\quad \Delta=a^2+H,\quad D=X+ac+\gamma H,\\
G=(q-1)(q-F)+\alpha,\quad R=G(q^2-1)+(M_C+qM_F)J,\\
L_0=EkY,\quad V=of-c,\quad U=jc-R,\quad I=ic^2,\\
A_0=\Delta(f^2-1).
\end{gather*}
The seven equations are
\begin{gather*}
 \tau^2-L_0(L_0+k)=1,\qquad D^2-\Delta c^2=1,\\
 A_0(V^2-y^2)+y^2=1,\qquad k-hE-R=1,\\
 (K_0+w)C=F+z(q-1),\qquad I^2=A_0,\qquad U=V.
\end{gather*}
The exact evaluator remains 59 operations.  By the inherited arithmetic
theorem, the native interface proved here, the unchanged Wang language,
and the literal K215 calibration, its positive unsolvability is provably
equivalent in ZFC to $\operatorname{Con}(\mathrm{ZFC})$.  The supplied
67-operation polynomial is rebound by the same complete tuple and keeps
its fifteen positive unknowns and exact degree 155.

\subsection{Reproduction and limits}
The builder reconstructs the literal tile language from all 430 pinned
machine actions without executing the machine.  The independent auditor
starts instead from the pinned tile list, imports no builder code,
reconstructs all four full integers, verifies the source Git blob identities,
and checks all native basis identities, exact carry and mask bounds, and
the prime controller.  It also verifies that both rebound arithmetic
sources are unchanged.  The complement lemma and inherited universal
compiler proofs, not finite diagnostic examples alone, supply the
unbounded mathematical conclusion.  No new end-to-end proof-assistant
certificate, consistency decision, or global optimum is claimed.


\section{Proof dependencies and reproducibility}
The archive contains the editable main report, a detailed Sun proof, universal proof notes for every new adapter, explicit division-free circuit JSON, materialized coefficient payloads, deterministic generators, independently written checkers, source hashes, and licensing notices.

The principal dependencies are separated as follows.
\begin{enumerate}
\item The literal K215 logical calibration and the native-window/Pell compiler identify the parent Diophantine predicate with a finite ZFC refutation. These results are imported from the pinned repository, not newly reconstructed in a proof assistant.
\item The exponent-set property has both a finite-field argument and exhaustive verification of all 52360 padded four-sums. The fixed-polynomial carry proof establishes both directions and the bounds needed before decoding.
\item The ten-integer route imports Sun's 2026 Lucas and bridge theorems. The nine-natural and thirteen/fourteen-integer routes instead import the existing Jones ratio theorems. The cofinal-base argument and the improved canonical-witness estimate are proved separately, rather than assumed from the statement of an existential theorem.
\item The twelve-positive and native complement constructions have their own two-sided local proofs. A modular evaluation test cannot replace their sign, integrality, carry, or completeness arguments.
\end{enumerate}
Run \texttt{python3 verify.py} from an extracted archive. The verifier authenticates the sealed input tree, copies it into a temporary directory, rebuilds the deterministic artifacts, runs each component's checks, and compares the rebuilt bytes with the release. The original archive contents are not modified. The final test ledger records the exact executed stages and rebuilt-file count. Python, SymPy and NumPy versions are recorded in \texttt{requirements.txt}; no Lean/Rocq installation is required for these finite checks.

The independent release review records both resolved issues and residual scope limits. In particular, an attempted elimination of the parent variable $w$ would require recovering $q\mid X$ from $q\mid XC$; the available coprimality argument is circular. That attempt is not counted as a thirty-variable quartic. Likewise, the earlier objection to squaring a slack variable does not by itself refute the proved restriction of a \emph{base} to scaled squares. The relevant distinction and source-level cofinality argument are preserved.

All results remain open to human scrutiny. The tests do not search for a solution of the consistency equation, prove unbounded unsolvability, establish consistency, or supply a new kernel certificate for the composite constructions.

\paragraph{Post-pin formalization work.}
At final review, commit \texttt{8747904f8b97} (full pin in the review source map) added the separate \texttt{Computability/DiophantineConsistency} work log. It reports no newly checked theorem and explicitly tracks computability, coefficient export, the formal P67 bridge, and object-language ZFC derivation as separate obligations. This is not a competing completed equation and does not change our frozen arithmetic baseline. We preserve its status file in the review directory. The present ordinary mathematical reductions do not claim to discharge those formalization obligations or export a syntactic ZFC derivation of the imported calibration. Their foundational claim is at the proof level of the explicitly imported compiler and logical-calibration results.


\begin{thebibliography}{9}
\bibitem{repo} V. Reshetnikov, \emph{ProveIt}, pinned snapshot cited above.\newline
\url{https://github.com/VladimirReshetnikov/ProveIt}
\bibitem{sun2021} Z.-W. Sun, \emph{Further results on Hilbert's Tenth Problem}, Science China Mathematics 64 (2021), 281--306.\newline
\url{https://arxiv.org/abs/1704.03504}
\bibitem{sun2026} Z.-W. Sun, \emph{Ten unknowns for Hilbert's tenth problem over the integers}, 2026 preprint, version 2.\newline
\url{https://arxiv.org/abs/2609.01594v2}
\end{thebibliography}
\end{document}
